\frfilename{mt262.tex}
\versiondate{11.8.15}
\copyrightdate{1995}

\def\dist{\mathop{\text{dist}}}

\def\chaptername{Perubahan variabel dalam integral}
\def\sectionname{Fungsi Lipschitz dan fungsi terdiferensialkan}

\newsection{262}

Sebagai persiapan untuk pekerjaan utama bab ini dalam \S263, saya
menyediakan satu bagian bagi dua kelas penting fungsi antar-ruang Euklides.
Yang sebenarnya kita perlukan ialah bahan yang pada hakikatnya elementer
sampai 262I, bersama lemma teknis 262M dan akibat-akibatnya.   Teorema
262Q tidak digunakan dalam jilid ini, walaupun saya yakin teorema itu
membuat pola-pola yang akan berkembang menjadi lebih alami dan mudah dipahami.

Seperti dalam \S261, $r$ (dan di sini juga $s$) akan selalu merupakan
bilangan bulat positif, dan `terukur', `terabaikan', `terintegralkan' akan
mengacu pada ukuran Lebesgue kecuali dinyatakan lain.

\leader{262A}{Fungsi Lipschitz} Misalkan
$\phi:D\to\BbbR^s$ suatu fungsi, dengan $D\subseteq\BbbR^r$.   
\cmmnt{Kita katakan bahwa }$\phi$ {\bf $\gamma$-Lipschitz}, dengan
$\gamma\in\coint{0,\infty}$, jika

\Centerline{$\|\phi(x)-\phi(y)\|\le \gamma \|x-y\|$}

\noindent untuk semua $x$, $y\in D$, dengan menuliskan
$\|x\|=\sqrt{\xi_1^2+\ldots+\xi_r^2}$ jika
$x=(\xi_1,\ldots,\xi_r)\in\BbbR^r$,
$\|z\|=\sqrt{\zeta_1^2+\ldots+\zeta_s^2}$ jika
$z=(\zeta_1,\ldots,\zeta_s)\in\BbbR^s$.   Dalam hal ini, $\gamma$ adalah
{\bf konstanta Lipschitz bagi $\phi$}.

Sebuah {\bf fungsi Lipschitz} adalah fungsi $\phi$ yang
$\gamma$-Lipschitz untuk suatu $\gamma\ge 0$.
\cmmnt{Perhatikan bahwa dalam hal ini $\phi$ mempunyai konstanta
Lipschitz terkecil (sebab jika $A$ adalah himpunan konstanta Lipschitz bagi
$\phi$, dan $\gamma_0=\inf A$, maka $\gamma_0$ merupakan konstanta Lipschitz
bagi $\phi$).}
\dvAnew{2015}Jelas bahwa fungsi Lipschitz kontinu (secara seragam).

\leader{262B}{}\cmmnt{ Kita memerlukan fakta-fakta mudah berikut.

\medskip

\noindent}{\bf Lemma} Misalkan $D\subseteq\BbbR^r$ suatu himpunan dan
$\phi:D\to\BbbR^s$ suatu fungsi.

(a)  $\phi$ Lipschitz jika dan hanya jika $\phi_i:D\to\Bbb R$
Lipschitz untuk setiap $i$, dengan menuliskan
$\phi(x)=(\phi_1(x),\ldots,\phi_s(x))$ untuk setiap $x\in
D=\dom\phi\subseteq\BbbR^r$.

(b) Dalam hal ini, terdapat fungsi Lipschitz
$\tilde\phi:\BbbR^r\to\BbbR^s$ yang memperluas $\phi$.

(c) Jika $r=s=1$ dan $D=[a,b]$ suatu interval,
maka $\phi$ Lipschitz jika dan hanya jika kontinu mutlak dan mempunyai
turunan terbatas.

\proof{{\bf (a)} Untuk sebarang $x$, $y\in D$ dan $i\le s$,

\Centerline{$|\phi_i(x)-\phi_i(y)|\le\|\phi(x)-\phi(y)\|
\le\sqrt{s}\sup_{j\le s}|\phi_j(x)-\phi_j(y)|$,}

\noindent sehingga setiap konstanta Lipschitz bagi $\phi$ merupakan
konstanta Lipschitz bagi setiap $\phi_i$, dan jika $\gamma_j$ adalah
konstanta Lipschitz bagi $\phi_j$ untuk setiap $j$, maka
$\sqrt{s}\sup_{j\le s}\gamma_j$ merupakan konstanta Lipschitz bagi $\phi$.

\medskip

{\bf (b)} Berdasarkan (a), cukup dipertimbangkan kasus $s=1$, sebab jika
setiap $\phi_i$ mempunyai perluasan Lipschitz $\tilde\phi_i$, kita dapat
menetapkan $\tilde\phi(x)=(\tilde\phi_1(x),\ldots,\tilde\phi_s(x))$ untuk
setiap $x$ guna memperoleh perluasan Lipschitz dari $\phi$.
Jadi ambil $s=1$; perhatikan bahwa kasus $D=\emptyset$ sepele;  maka
andaikan $D\ne\emptyset$.   Misalkan $\gamma$ suatu konstanta Lipschitz
bagi $\phi$, dan tuliskan

\Centerline{$\tilde\phi(z)=\sup_{y\in D}\phi(y)-\gamma\|y-z\|$}

\noindent untuk setiap $z\in\BbbR^r$.   Jika $x\in D$, maka untuk sebarang
$z\in\BbbR^r$ dan $y\in D$,

\Centerline{$\phi(y)-\gamma\|y-z\|
\le\phi(x)+\gamma\|y-x\|-\gamma\|y-z\|
\le\phi(x)+\gamma\|z-x\|$,}

\noindent sehingga $\tilde\phi(z)\le\phi(x)+\gamma\|z-x\|$;  khususnya,
ini menunjukkan bahwa $\tilde\phi(z)<\infty$.   Selain itu, jika $z\in D$,
harus berlaku

\Centerline{$\phi(z)-\gamma\|z-z\|\le\tilde\phi(z)
\le\phi(z)+\gamma\|z-z\|$,}

\noindent sehingga $\tilde\phi$ memperluas $\phi$.   Akhirnya, jika $w$,
$z\in\BbbR^r$ dan $y\in D$,

\Centerline{$\phi(y)-\gamma\|y-w\|
\le\phi(y)-\gamma\|y-z\|+\gamma\|w-z\|
\le\tilde\phi(z)+\gamma\|w-z\|$;}

\noindent dan dengan mengambil supremum atas $y\in D$,

\Centerline{$\tilde\phi(w)\le\tilde\phi(z)+\gamma\|w-z\|$.}

\noindent Karena $w$ dan $z$ sebarang, $\tilde\phi$ Lipschitz.

\medskip

{\bf (c)(i)} Andaikan $\phi$ $\gamma$-Lipschitz.  Jika $\epsilon>0$ dan
$a\le a_1\le b_1\le\ldots\le a_n\le b_n\le b$ serta
$\sum_{i=1}^n(b_i-a_i)\le\epsilon/(1+\gamma)$, maka

\Centerline{$\sum_{i=1}^n|\phi(b_i)-\phi(a_i)|
\le\sum_{i=1}^n\gamma|b_i-a_i|
\le\epsilon$.}

\noindent Karena $\epsilon$ sebarang, $\phi$ kontinu mutlak.
Jika $x\in[a,b]$ dan $\phi'(x)$ terdefinisi, maka

\Centerline{$|\phi'(x)|
=\lim_{y\to x}\Bover{|\phi(y)-\phi(x)|}{|y-x|}\le\gamma$,}

\noindent sehingga $\phi'$ terbatas.

\medskip

\quad{\bf (ii)} Sekarang andaikan $\phi$ kontinu mutlak dan
$|\phi'(x)|\le\gamma$ untuk setiap $x\in\dom\phi'$, dengan
$\gamma\ge 0$.   Maka setiap kali $a\le x\le y\le b$,

\Centerline{$|\phi(y)-\phi(x)|=|\int_x^y\phi'|\le\int_x^y|\phi'|
\le\gamma(y-x)$}

\noindent (menggunakan 225E untuk kesamaan pertama).   Karena $x$ dan $y$
sebarang, $\phi$ $\gamma$-Lipschitz.
}%end of proof of 262B

\cmmnt{
\leader{262C}{Catatan} Argumen untuk (b) di atas menunjukkan bahwa jika
$\phi:D\to\Bbb R$ suatu fungsi Lipschitz, dengan
$D\subseteq\Bbb R^r$, maka $\phi$ mempunyai perluasan ke $\BbbR^r$ dengan
konstanta Lipschitz yang sama.   Bahkan, jika $\phi:D\to\BbbR^s$ suatu
fungsi Lipschitz, maka $\phi$ mempunyai perluasan
$\tilde\phi:\BbbR^r\to\BbbR^s$ dengan konstanta Lipschitz yang sama;
ini adalah `teorema Kirszbraun' ({\smc Kirszbraun 34}, atau
{\smc Federer 69}, 2.10.43).
}%end of comment

\leader{262D}{Proposisi} Jika $\phi:D\to\BbbR^r$ suatu fungsi
$\gamma$-Lipschitz, dengan $D\subseteq\BbbR^r$, maka
$\mu^*\phi[A]\le\gamma^r\mu^*A$ untuk setiap $A\subseteq D$, dengan $\mu$
ukuran Lebesgue pada $\BbbR^r$.   Khususnya,
$\phi[D\cap A]$ terabaikan untuk setiap himpunan terabaikan
$A\subseteq\BbbR^r$.

\proof{ Misalkan $\epsilon>0$.   Berdasarkan 261F, terdapat barisan bola
tertutup $\sequencen{B_n}=\sequencen{B(x_n,\delta_n)}$ di
$\Bbb R^r$ yang menutupi $A$, sedemikian sehingga
$\sum_{n=0}^{\infty}\mu B_n\le\mu^* A+\epsilon$ dan
$\sum_{n\in\Bbb N\setminus K}\mu B_n\le\epsilon$, dengan
$K=\{n:n\in\Bbb N,\,x_n\in A\}$.   Tetapkan

\Centerline{$L=\{n:n\in\Bbb N\setminus K,\,B_n\cap D\ne\emptyset\}$,}

\noindent dan untuk $n\in L$ pilih $y_n\in D\cap B_n$.   Sekarang tetapkan

$$\eqalign{B'_n
&=B(\phi(x_n),\gamma\delta_n)\text{ jika }n\in K,\cr
&=B(\phi(y_n),2\gamma\delta_n)\text{ jika }n\in L,\cr
&=\emptyset\text{ jika }n\in\Bbb N\setminus(K\cup L).\cr}$$

\noindent Maka $\phi[B_n\cap D]\subseteq B'_n$ untuk setiap $n$, sehingga
$\phi[D\cap A]\subseteq\bigcup_{n\in\Bbb N}B'_n$, dan

$$\eqalign{\mu^*\phi[A\cap D]
&\le\sum_{n=0}^{\infty}\mu B'_n
=\gamma^r\sum_{n\in K}\mu B_n+2^r\gamma^r\sum_{n\in L}\mu B_n\cr
&\le\gamma^r(\mu^*A+\epsilon)+2^r\gamma^r\epsilon.\cr}$$

\noindent Karena $\epsilon$ sebarang, $\mu^*\phi[A\cap
D]\le\gamma^r\mu^*A$, sebagaimana diklaim.
}%end of proof of 262D

\vleader{72pt}{262E}{Akibat} Misalkan $\phi:D\to\BbbR^r$ suatu fungsi
Lipschitz injektif, dengan $D\subseteq\BbbR^r$, dan $f$ suatu fungsi
terukur dari suatu himpunan bagian $\BbbR^r$ ke $\Bbb R$.

(a) Jika $\phi^{-1}$ terdefinisi hampir di mana-mana dalam suatu himpunan $H$
dari $\BbbR^r$ dan $f$ terdefinisi hampir di mana-mana dalam
$\BbbR^r$, maka $f\phi^{-1}$ terdefinisi hampir di mana-mana dalam $H$.

(b) Jika $E\subseteq D$ terukur Lebesgue, maka $\phi[E]$ terukur.

(c) Jika $D$ terukur, maka $f\phi^{-1}$ terukur.

\proof{ Tetapkan

\Centerline{$C=\dom(f\phi^{-1})
=\{y:y\in\phi[D],\,\phi^{-1}(y)\in\dom f\}=\phi[D\cap\dom f]$.}

\medskip

{\bf (a)} Karena $f$ terdefinisi hampir di mana-mana,
$\phi[D\setminus\dom f]$ terabaikan.   Tetapi sekarang

\Centerline{$C=\phi[D]\setminus \phi[D\setminus\dom f]
=\dom\phi^{-1}\setminus\phi[D\setminus\dom f]$,}

\noindent sehingga

\Centerline{$H\setminus C
\subseteq(H\setminus\dom\phi^{-1})\cup\phi[D\setminus\dom f]$}

\noindent terabaikan.

\medskip

{\bf (b)} Sekarang andaikan $E\subseteq D$ dan $E$ terukur.
Misalkan $\sequencen{F_n}$ suatu barisan himpunan bagian tertutup dan
terbatas dari $E$ sedemikian sehingga
$\mu(E\setminus\bigcup_{n\in\Bbb N}F_n)=0$ (134Fb).   Karena $\phi$
Lipschitz, fungsi ini kontinu, sehingga $\phi[F_n]$ kompak, jadi tertutup,
dan karena itu terukur untuk setiap $n$ (2A2E, 115G); selain itu,
$\phi[E\setminus\bigcup_{n\in\Bbb N}F_n]$ terabaikan berdasarkan 262D,
dan karena itu terukur.   Jadi

\Centerline{$\phi[E]
=\phi[E\setminus\bigcup_{n\in\Bbb N}F_n]
  \cup\bigcup_{n\in\Bbb N}\phi[F_n]
$}

\noindent terukur.

\medskip

{\bf (c)} Untuk sebarang $a\in\Bbb R$, ambil himpunan terukur
$E\subseteq\BbbR^r$ sedemikian sehingga
$\{x:f(x)\ge a\}=E\cap\dom f$.   Maka

\Centerline{$\{y:y\in C,\,f\phi^{-1}(y)\ge a\}
=C\cap\phi[D\cap E]$.}

\noindent Tetapi $\phi[D\cap E]$ terukur berdasarkan (b), sehingga
$\{y:f\phi^{-1}(y)\ge a\}$ terukur relatif dalam $C$.   Karena $a$
sebarang, $f\phi^{-1}$ terukur.
}%end of proof of 262E

\leader{262F}{\dvrocolon{Keterdiferensialan}}\cmmnt{ Sekarang saya
beralih ke kelas fungsi yang sifat-sifatnya akan mengisi sebagian besar
sisa bab ini.

\medskip

\noindent}{\bf Definisi}  Misalkan $\phi$ suatu fungsi dari himpunan
bagian $D=\dom\phi$ dari $\BbbR^r$ ke $\BbbR^s$.

\spheader 262Fa $\phi$ {\bf terdiferensialkan} di $x\in D$ jika terdapat
matriks real $s\times r$, yaitu $T$, sedemikian sehingga

\Centerline{$\lim_{y\to x}
\Bover{\|\phi(y)-\phi(x)-T(y-x)\|}{\|y-x\|}=0$;}

\noindent dalam hal ini kita dapat menuliskan $T=\phi'(x)$.

\spheader 262Fb Saya akan mengatakan bahwa $\phi$ {\bf terdiferensialkan
relatif terhadap domainnya} di $x$, dan bahwa $T$ adalah {\bf suatu}
turunan dari $\phi$ di $x$, jika $x\in D$ dan untuk setiap $\epsilon>0$
terdapat $\delta>0$ sedemikian sehingga
$\|\phi(y)-\phi(x)-T(y-x)\|\le\epsilon\|y-x\|$ untuk setiap
$y\in B(x,\delta)\cap D$.

\cmmnt{
\vleader{72pt}{262G}{Catatan (a)} Definisi baku dalam 262Fa, yang
melibatkan limit dari semua arah
`$\lim_{y\to x}$', secara tersirat mensyaratkan $\phi$ terdefinisi pada
suatu bola tak-sepele yang berpusat di $x$, agar kita dapat menghitung
$\phi(y)-\phi(x)-T(y-x)$ untuk semua $y$ yang cukup dekat dengan $x$.
Keuntungannya ialah turunan $T=\phi'(x)$ terdefinisi secara tunggal
(sebab jika $\lim_{z\to \tbf{0}}\Bover{\|T_1z-T_2z\|}{\|z\|}=0$ maka

\Centerline{$\Bover{\|(T_1-T_2)z\|}{\|z\|}=\lim_{\alpha\to
0}\Bover{\|T_1(\alpha
z)-T_2(\alpha z)\|}{\|\alpha z\|}=0$}

\noindent untuk setiap $z$ taknol, sehingga $T_1-T_2$ harus merupakan
matriks nol). Untuk keperluan kita di sini, ada manfaatnya melonggarkan
hal ini sedikit menjadi bentuk dalam 262Fb, agar kita tidak perlu memberi
perhatian khusus pada batas $\dom\phi$.   Khususnya, kita mendapati bahwa
jika $T$ suatu turunan dari $\phi:D\to\BbbR^s$ relatif terhadap domainnya
di $x$, dan $x\in D'\subseteq D$, maka $T$ merupakan turunan dari
$\phi\restr D'$ relatif terhadap domainnya di $x$.

\header{262Gb}{\bf (b)} Jika Anda belum pernah menjumpai konsep
`keterdiferensialan' ini, tetapi cukup mengenal pendiferensialan parsial,
perlu ditekankan bahwa konsep fungsi `terdiferensialkan' (setidaknya dalam
pengertian ketat yang dituntut 262Fa) benar-benar lebih kuat daripada
konsep fungsi `terdiferensialkan secara parsial'. Untuk keperluan
perhitungan, metode paling berguna untuk menemukan turunan sejati diberikan
oleh 262Id di bawah.   Sebagai contoh sederhana fungsi yang mempunyai
semua turunan parsial tetapi tidak terdiferensialkan di setiap titik,
perhatikan $\phi:\BbbR^2\to\Bbb R$ yang didefinisikan oleh

$$\eqalign{\phi(\xi_1,\xi_2)
&={{\xi_1\xi_2}\over{\xi_1^2+\xi_2^2}}
 \text{ jika }\xi_1^2+\xi_2^2\ne 0,\cr
&=0\text{ jika }\xi_1=\xi_2=0.\cr}$$

\noindent Maka $\phi$ bahkan tidak kontinu di $\tbf{0}$, walaupun kedua
turunan parsial $\pd{\phi}{\xi_j}$ terdefinisi di mana-mana.

\header{262Gc}{\bf (c)} Dalam definisi di atas, saya menyebut turunan
sebagai suatu matriks.   Secara tepat, turunan fungsi yang didefinisikan
pada suatu himpunan bagian $\BbbR^r$ dan bernilai dalam $\BbbR^s$
seharusnya dipandang sebagai operator linear terbatas dari $\BbbR^r$ ke
$\BbbR^s$;  perumusan dengan matriks dapat diterima karena terdapat
korespondensi satu-satu alami antara matriks real $s\times r$ dan operator
linear dari $\BbbR^r$ ke $\BbbR^s$, dan semua operator linear ini terbatas.
Saya menggunakan deskripsi `matriks' karena deskripsi ini membuat beberapa
perhitungan lebih langsung; khususnya, hubungan antara $\phi'$ dan turunan
parsial $\phi$ (262Ic), serta gagasan determinan $\det\phi'(x)$, yang
digunakan sepanjang \S\S263 dan 265.
}%end of comment

\leader{262H}{Norma matriks} Beberapa perhitungan di bawah akan bertumpu
pada gagasan `norma' matriks.   Norma yang akan saya gunakan\cmmnt{
(sebenarnya, untuk keperluan kita di sini, norma apa pun memadai)} ialah
`norma operator', yang didefinisikan dengan menetapkan

\Centerline{$\|T\|=\sup\{\|Tx\|:x\in\BbbR^r,\,\|x\|\le 1\}$}

\noindent untuk sebarang matriks $s\times r$, yaitu $T$.   \cmmnt{Untuk
fakta-fakta dasar mengenai norma-norma ini, lihat 2A4F-2A4G.   Hal berikut
juga akan berguna.}

\header{262Ha}{\bf (a)} Jika\cmmnt{ semua koefisien $T$ kecil, maka
$\|T\|$ juga kecil;  sesungguhnya, jika}
$T=\langle\tau_{ij}\rangle_{i\le s,j\le r}$\cmmnt{, dan
$\|x\|\le 1$,} maka
\cmmnt{$|\xi_j|\le 1$ untuk setiap $j$, sehingga

\Centerline{$\|Tx\|
=\bigl(\sum_{i=1}^s(\sum_{j=1}^r\tau_{ij}\xi_j)^2\bigr)^{1/2}
\le \bigl(\sum_{i=1}^s(\sum_{j=1}^r|\tau_{ij}|)^2\bigr)^{1/2}
\le r\sqrt{s}\max_{i\le s,j\le r}|\tau_{ij}|$,}

\noindent dan} $\|T\|\le r\sqrt{s}\max_{i\le s,j\le r}|\tau_{ij}|$.
\cmmnt{(Ini ketaksamaan yang sangat kasar.   Ketaksamaan yang lebih baik
ada dalam 262Yb.   Namun khususnya ini memberi tahu kita bahwa $\|T\|$
selalu berhingga.)
}

\header{262Hb}{\bf (b)}\dvro{ $|\tau_{ij}|\le\|T\|$ untuk semua $i$, $j$.}
{ Jika $\|T\|$ kecil, maka semua koefisien $T$ juga kecil; sesungguhnya,
dengan menuliskan $e_j$ untuk vektor satuan ke-$j$ dari $\BbbR^r$, maka
koordinat ke-$i$ dari $Te_j$ adalah $\tau_{ij}$, sehingga
$|\tau_{ij}|\le\|Te_j\|\le\|T\|$.}

\leader{262I}{Lemma} Misalkan $\phi:D\to\BbbR^s$ suatu fungsi, dengan
$D\subseteq\BbbR^r$.   Untuk $i\le s$, misalkan $\phi_i:D\to\Bbb R$
koordinat ke-$i$-nya, sehingga $\phi(x)=(\phi_1(x),\ldots,\phi_s(x))$
untuk $x\in D$.

(a) Jika $\phi$ terdiferensialkan relatif terhadap domainnya di $x\in D$,
maka $\phi$ kontinu di $x$.

(b) Jika $x\in D$, maka $\phi$ terdiferensialkan relatif terhadap
domainnya di $x$ jika dan hanya jika setiap $\phi_i$ terdiferensialkan
relatif terhadap domainnya di $x$.

(c) Jika $\phi$ terdiferensialkan di $x\in D$, maka semua turunan parsial
$\pd{\phi_i}{\xi_j}$ dari $\phi$ terdefinisi di $x$, dan turunan $\phi$
di $x$ adalah matriks
$\langle\pd{\phi_i}{\xi_j}(x)\rangle_{i\le s,j\le r}$.

(d) Jika semua turunan parsial $\pd{\phi_i}{\xi_j}$, untuk $i\le s$ dan
$j\le r$, terdefinisi di suatu lingkungan dari $x\in D$ dan kontinu di $x$,
maka $\phi$ terdiferensialkan di $x$.

\proof{{\bf (a)} Misalkan $T$ suatu turunan dari $\phi$ di $x$. Dengan
menerapkan definisi 262Fb dengan $\epsilon=1$, kita melihat bahwa terdapat
$\delta>0$ sedemikian sehingga
\Centerline{$\|\phi(y)-\phi(x)-T(y-x)\|\le\|y-x\|$}

\noindent setiap kali $y\in D$ dan $\|y-x\|\le\delta$.   Sekarang

\Centerline{$\|\phi(y)-\phi(x)\|\le\|T(y-x)\|+\|y-x\|
\le(1+\|T\|)\|y-x\|$}

\noindent setiap kali $y\in D$ dan $\|y-x\|\le\delta$, sehingga $\phi$
kontinu di $x$.

\medskip

{\bf (b)(i)} Jika $\phi$ terdiferensialkan relatif terhadap domainnya di
$x\in D$, misalkan $T$ suatu turunan dari $\phi$ di $x$. Untuk $i\le s$,
misalkan $T_i$ matriks $1\times r$ yang terdiri atas baris ke-$i$ dari
$T$. Misalkan $\epsilon>0$. Maka terdapat $\delta>0$ sedemikian sehingga

$$\eqalign{|\phi_i(y)-\phi_i(x)-T_i(y-x)|
&\le\|\phi(y)-\phi(x)-T(y-x)\|\cr
&\le\epsilon\|y-x\|\cr}$$

\noindent setiap kali $y\in D$ dan $\|y-x\|\le\delta$, sehingga $T_i$
merupakan turunan dari $\phi_i$ di $x$.

\medskip

\quad{\bf (ii)} Jika setiap $\phi_i$ terdiferensialkan relatif terhadap
domainnya di $x$, dengan turunan yang bersesuaian $T_i$, misalkan $T$
matriks $s\times r$ dengan baris $T_1,\ldots,T_s$. Untuk $\epsilon>0$
yang diberikan, bagi setiap $i\le s$ terdapat $\delta_i>0$ sedemikian
sehingga

\Centerline{$|\phi_i(y)-\phi_i(x)-T_i(y-x)|\le\epsilon\|y-x\|$ setiap kali
$y\in D$, $\|y-x\|\le\delta_i$;}

\noindent tetapkan $\delta=\min_{i\le s}\delta_i>0$; maka jika $y\in D$
dan $\|y-x\|\le\delta$, akan berlaku

\Centerline{$\|\phi(y)-\phi(x)-T(y-x)\|^2
=\sum_{i=1}^s|\phi_i(y)-\phi_i(x)-T_i(y-x)|^2
\le s\epsilon^2\|y-x\|^2$,}

\noindent sehingga

\Centerline{$\|\phi(y)-\phi(x)-T(y-x)\|
\le\epsilon\sqrt{s}\|y-x\|$.}

\noindent Karena $\epsilon$ sebarang, $T$ merupakan turunan dari $\phi$
di $x$.

\medskip

{\bf (c)} Tetapkan $T=\phi'(x)$.   Kita mempunyai

\Centerline{$\lim_{y\to x}\Bover{\|\phi(y)-\phi(x)-T(y-x)\|}{\|y-x\|}
=0$;}

\noindent tetapkan $j\le r$, dan perhatikan $y=x+\eta e_j$, dengan
$e_j=(0,\ldots,0,1,0,\ldots,0)$ vektor satuan ke-$j$ dalam $\BbbR^r$.
Maka harus berlaku

\Centerline{$\lim_{\eta\to 0}\Bover{\|\phi(x+\eta e_j)-\phi(x)-\eta
T(e_j)\|}{|\eta|}
=0$.}

\noindent Dengan memperhatikan koordinat ke-$i$ dari
$\phi(x+\eta e_j)-\phi(x)-\eta T(e_j)$, kita memperoleh


\Centerline{$|\phi_i(x+\eta e_j)-\phi_i(x)-\tau_{ij}\eta|
\le\|\phi(x+\eta e_j)-\phi(x)-\eta T(e_j)\|$,}

\noindent dengan $\tau_{ij}$ koefisien ke-$(i,j)$ dari $T$; sehingga

\Centerline{$\lim_{\eta\to 0}
\Bover{|\phi_i(x+\eta e_j)-\phi_i(x)-\tau_{ij}\eta|}{|\eta|}
=0$.}

\noindent Tetapi ini tepat mengatakan bahwa turunan parsial
$\pd{\phi_i}{\xi_j}(x)$ ada dan sama dengan $\tau_{ij}$, sebagaimana
dinyatakan.

\medskip

{\bf (d)} Sekarang andaikan turunan parsial $\pd{\phi_i}{\xi_j}$
terdefinisi di dekat $x$ dan kontinu di $x$. Misalkan $\epsilon>0$.
Misalkan $\delta>0$ sedemikian sehingga

\Centerline{$|\pd{\phi_i}{\xi_j}(y)-\tau_{ij}|\le\epsilon$}

\noindent setiap kali $\|y-x\|\le\delta$, dengan menuliskan
$\tau_{ij}=\pd{\phi_i}{\xi_j}(x)$.   Sekarang andaikan
$\|y-x\|\le\delta$.   Tetapkan

\Centerline{$y=(\eta_1,\ldots,\eta_r)$,
\quad$x=(\xi_1,\ldots,\xi_r)$,}

\Centerline{$y_j=(\eta_1,\ldots,\eta_j,\xi_{j+1},\ldots,\xi_r)$ untuk
$0\le
j\le r$,}

\noindent sehingga $y_0=x$, $y_r=y$, dan ruas garis antara $y_{j-1}$ dan
$y_{j}$ seluruhnya berada dalam jarak $\delta$ dari $x$ setiap kali
$1\le j\le r$, sebab jika $z=(\zeta_1,\ldots,\zeta_r)$ terletak pada
ruas ini, maka $\zeta_i$ berada di antara $\xi_i$ dan $\eta_i$ untuk
setiap $i$. Berdasarkan teorema nilai rata-rata biasa bagi fungsi real
terdiferensialkan, yang diterapkan pada fungsi

\Centerline{$t\mapsto
\phi_i(\eta_1,\ldots,\eta_{j-1},t,\xi_{j+1},\ldots,\xi_r)$,}

\noindent untuk setiap $i\le s$, $j\le r$ terdapat titik $z_{ij}$ pada
ruas garis antara $y_{j-1}$ dan $y_j$ sedemikian sehingga

\Centerline{$\phi_i(y_j)-\phi_i(y_{j-1})
=(\eta_j-\xi_j)\pd{\phi_i}{\xi_j}(z_{ij})$.}

\noindent Tetapi

\Centerline{$|\pd{\phi_i}{\xi_j}(z_{ij})-\tau_{ij}|\le\epsilon$,}
\noindent sehingga

\Centerline{$|\phi_i(y_j)-\phi_i(y_{j-1})
-\tau_{ij}(\eta_j-\xi_j)|\le\epsilon|\eta_j-\xi_j|\le\epsilon\|y-x\|$.}

\noindent Dengan menjumlahkan atas $j$,

\Centerline{$|\phi_i(y)-\phi_i(x)-\sum_{j=1}^r\tau_{ij}(\eta_j-\xi_j)|
\le r\epsilon\|y-x\|$}

\noindent untuk setiap $i$.   Dengan menjumlahkan kuadrat-kuadratnya dan
mengambil akar kuadrat,

\Centerline{$\|\phi(y)-\phi(x)-T(y-x)\|\le \epsilon r\sqrt{s}\|y-x\|$,}

\noindent dengan $T=\langle\tau_{ij}\rangle_{i\le s,j\le r}$. Ini berlaku
setiap kali $\|y-x\|\le\delta$.   Karena $\epsilon$ sebarang,
$\phi'(x)=T$ terdefinisi.
}%end of proof of 262I

\cmmnt{
\leader{262J}{Catatan} Saya tidak yakin apakah saya perlu meminta maaf
atas notasi $\pd{}{\xi_j}$.   Dalam rumus seperti
$(\eta_j-\xi_j)\pd{\phi_i}{\xi_j}(z_{ij})$ di atas, kedua kemunculan
$\xi_j$ sangat mudah membingungkan.   Namun saya rasa orang yang
beriktikad baik tidak akan tersesat, asalkan label $\xi_j$ (atau simbol
apa pun yang digunakan untuk mewakili variabel terkait) dijelaskan secara
memadai ketika domain $\phi$ pertama kali diperkenalkan (dan selalu diingat
bahwa dalam pendiferensialan parsial, kita tidak hanya menggerakkan satu
variabel -- suatu $\xi_j$ dalam konteks sekarang -- tetapi juga menahan
tetap daftar variabel lain yang tidak dicantumkan dalam notasi). Saya yakin
notasi tradisional $\pd{}{\xi_j}$ bertahan karena alasan yang kuat, dan
saya ingin menyambut mereka yang lebih nyaman dengannya daripada dengan
berbagai alternatif yang pernah diusulkan tetapi tidak pernah berakar.
}%end of comment

\cmmnt{
\leader{262K}{Meninjau kembali fungsi Cantor} Ada manfaatnya memeriksa
kembali contoh-contoh 134H-134I dalam terang pertimbangan sekarang.
Misalkan $f:[0,1]\to[0,1]$ fungsi Cantor (134H) dan tetapkan
$g(x)=\bover12(x+f(x))$ untuk $x\in [0,1]$.   Maka
$g:[0,1]\to[0,1]$ suatu homeomorfisme (134I); tetapkan
$\phi=g^{-1}:[0,1]\to[0,1]$.   Kita melihat bahwa jika
$0\le x\le y\le 1$ maka $g(y)-g(x)\ge\bover12(y-x)$; secara ekuivalen,
$\phi(y)-\phi(x)\le 2(y-x)$ setiap kali $0\le x\le y\le 1$, sehingga
$\phi$ merupakan fungsi Lipschitz, dan karena itu kontinu mutlak (262Bc).
Jika $D=\{x:\phi'(x)$ terdefinisi$\}$, maka $[0,1]\setminus D$ terabaikan
(225Cb), sehingga $[0,1]\setminus\phi[D]=\phi[\,[0,1]\setminus D]$
terabaikan (262Da). Saya mencatat dalam 134I bahwa terdapat fungsi terukur
$h:[0,1]\to\Bbb R$ sedemikian sehingga komposisi $h\phi$ tidak terukur;
sekarang $h(\phi\restr D)=(h\phi)\restr D$ tidak mungkin terukur, meskipun
$\phi\restr D$ terdiferensialkan.
}%end of comment

\leader{262L}{}\cmmnt{ Akan berguna jika kita dapat memakai hasil langsung
berikut.

\medskip

\noindent}{\bf Lemma} Andaikan $D\subseteq\BbbR^r$ dan $x\in\BbbR^r$
sedemikian sehingga
$\lim_{\delta\downarrow 0}
 \Bover{\mu^*(D\cap B(x,\delta))}{\mu B(x,\delta)}=1$.   Maka
$\lim_{z\to\tbf{0}}\Bover{\rho(x+z,D)}{\|z\|}=0$, dengan
$\rho(x+z,D)=\inf_{y\in D}\|x+z-y\|$.

\proof{ Misalkan $\epsilon>0$. Misalkan $\delta_0>0$ sedemikian sehingga

\Centerline{$\mu^*(D\cap B(x,\delta))
>(1-(\Bover{\epsilon}{1+\epsilon})^r)\mu B(x,\delta)$}

\noindent setiap kali $0<\delta\le\delta_0$. Ambil sebarang $z$ sedemikian
sehingga $0<\|z\|\le\delta_0/(1+\epsilon)$.   \Quer\ Andaikan, jika
mungkin, bahwa $\rho(x+z,D)>\epsilon\|z\|$. Maka
$B(x+z,\epsilon\|z\|)\subseteq B(x,(1+\epsilon)\|z\|)\setminus D$,
sehingga

$$\eqalign{\mu^*(D\cap B(x,(1+\epsilon)\|z\|))
&\le\mu B(x,(1+\epsilon)\|z\|)-\mu B(x+z,\epsilon\|z\|)\cr
&=(1-(\Bover{\epsilon}{1+\epsilon})^r)\mu B(x,(1+\epsilon)\|z\|),\cr}$$

\noindent yang mustahil, sebab $(1+\epsilon)\|z\|\le\delta_0$.\
\BanG\ Jadi $\rho(x+z,D)\le\epsilon\|z\|$. Karena $\epsilon$ sebarang,
ini membuktikan hasil tersebut.
}%end of proof of 262L

\cmmnt{\medskip

\noindent{\bf Catatan} Ada istilah untuk hal ini; lihat 261Yg.}

\leader{262M}{}\cmmnt{ Sekarang saya sampai pada hasil pertama yang
menghubungkan fungsi Lipschitz dengan fungsi terdiferensialkan. Saya
mendekatinya melalui suatu lemma penting yang akan menjadi landasan \S263.

\medskip

\noindent}{\bf Lemma} Misalkan $\phi$ suatu fungsi dari himpunan bagian
$D$ dari $\BbbR^r$ ke $\BbbR^s$ yang terdiferensialkan di setiap titik
domainnya. Untuk setiap $x\in D$, misalkan $T(x)$ suatu turunan dari
$\phi$. Misalkan $M_{sr}$ himpunan matriks $s\times r$ dan
$\zeta:A\to\ooint{0,\infty}$ suatu fungsi positif ketat, dengan
$A\subseteq M_{sr}$ suatu himpunan takkosong yang memuat $T(x)$ untuk
setiap $x\in D$. Maka kita dapat menemukan barisan
$\sequencen{D_n}$, $\sequencen{T_n}$ sedemikian sehingga

(i) $\sequencen{D_n}$ adalah partisi $D$ menjadi himpunan-himpunan yang
terukur relatif dalam $D$\cmmnt{, yaitu irisan $D$ dengan himpunan bagian
terukur dari $\BbbR^r$};

(ii) $T_n\in A$ untuk setiap $n$;

(iii) $\|\phi(x)-\phi(y)-T_n(x-y)\|\le\zeta(T_n)\|x-y\|$ untuk setiap
$n\in\Bbb N$ dan $x$, $y\in D_n$;

(iv) $\|T(x)-T_n\|\le\zeta(T_n)$ untuk setiap $x\in D_n$.

\proof{{\bf (a)} Langkah pertama ialah mencatat bahwa terdapat barisan
$\sequencen{S_n}$ dalam $A$ sedemikian sehingga

\Centerline{$A
\subseteq\bigcup_{n\in\Bbb N}\{T:T\in M_{sr},\,\|T-S_n\|<\zeta(S_n)\}$.}

\noindent\Prf\ (Tentu saja ini merupakan hasil baku mengenai ruang metrik
terpisahkan.) Tuliskan $Q$ untuk himpunan matriks dalam $M_{sr}$ dengan
koefisien rasional; maka terdapat bijeksi alami antara $Q$ dan
$\Bbb Q^{sr}$, sehingga $Q$ dan $Q\times\Bbb N$ terhitung. Daftarkan
$Q\times\Bbb N$ sebagai $\sequencen{(R_n,k_n)}$. Untuk setiap
$n\in\Bbb N$, pilih $S_n\in A$ menurut aturan

\quad --- jika terdapat $S\in A$ sedemikian sehingga
$\{T:\|T-R_n\|\le 2^{-k_n}\}\subseteq\{T:\|T-S\|<\zeta(S)\}$,
ambil $S$ seperti itu sebagai $S_n$;

\quad --- jika tidak, ambil $S_n$ sebagai sebarang anggota $A$.

\noindent Saya nyatakan bahwa ini berhasil. Misalkan $S\in A$. Maka
$\zeta(S)>0$; ambil $k\in\Bbb N$ sedemikian sehingga
$2^{-k}<\zeta(S)$. Ambil $R^*\in Q$ sedemikian sehingga
$\|R^*-S\|<\min(\zeta(S)-2^{-k},2^{-k})$; ini mungkin karena $\|R-S\|$
akan kecil setiap kali semua koefisien $R$ cukup dekat dengan koefisien
$S$ yang bersesuaian (262Ha), dan kita dapat memilih bilangan rasional
untuk mencapainya. Misalkan $n\in\Bbb N$ sedemikian sehingga $R^*=R_n$
dan $k=k_n$. Maka

\Centerline{$\{T:\|T-R_n\|\le 2^{-k_n}\}
\subseteq\{T:\|T-S\|<\zeta(S)\}$}

\noindent (sebab $\|T-S\|\le\|T-R_n\|+\|R_n-S\|$), sehingga kita pasti
memilih $S_n$ menurut bagian pertama aturan di atas, dan

\Centerline{$S\in\{T:\|T-R_n\|\le 2^{-k_n}\}
\subseteq\{T:\|T-S_n\|<\zeta(S_n)\}$.}

\noindent Karena $S$ sebarang, ini membuktikan hasil tersebut.\ \Qed

\medskip

{\bf (b)} Daftarkan $\Bbb Q^r\times\Bbb Q^r\times\Bbb N$ sebagai
$\sequencen{(q_n,q'_n,m_n)}$. Untuk setiap $n\in\Bbb N$, tetapkan

$$\eqalign{H_n&=\{x:x\in[q_n,q'_n]\cap D,\,
\|\phi(y)-\phi(x)-S_{m_n}(y-x)\|\le\zeta(S_{m_n})\|y-x\|\cr
&\mskip370mu
    \text{untuk setiap }y\in[q_n,q'_n]\cap D\}\cr
&=[q_n,q'_n]\cap D\cap\bigcap_{y\in[q_n,q'_n]\cap D}
       \{x:x\in D,\cr
&\mskip270mu
    \|\phi(y)-\phi(x)-S_{m_n}(y-x)\|\le\zeta(S_{m_n})\|y-x\|\}.\cr}$$

\noindent Karena $\phi$ kontinu, $H_n=D\cap\overline{H}_n$, dengan
$\overline{H}_n$ menyatakan penutupan $H_n$, sehingga $H_n$ terukur
relatif dalam $D$. Perhatikan bahwa jika $x$, $y\in H_n$, maka
$y\in D\cap[q_n,q'_n]$, sehingga

\Centerline{$\|\phi(y)-\phi(x)-S_{m_n}(y-x)\|\le\zeta(S_{m_n})\|y-x\|$.}


Tetapkan

\Centerline{$H_n'=\{x:x\in H_n,\,\|T(x)-S_{m_n}\|\le\zeta(S_{m_n})\}$.}

\medskip

{\bf (c)} $D=\bigcup_{n\in\Bbb N}H'_n$.   \Prf\ Misalkan $x\in D$.
Maka $T(x)\in A$, sehingga terdapat $k\in\Bbb N$ sedemikian sehingga
$\|T(x)-S_k\|<\zeta(S_k)$. Misalkan $\delta>0$ sedemikian sehingga

\Centerline{$\|\phi(y)-\phi(x)-T(x)(y-x)\|
\le(\zeta(S_k)-\|T(x)-S_k\|)\|x-y\|$}

\noindent setiap kali $y\in D$ dan $\|y-x\|\le\delta$.   Maka

$$\eqalign{\|\phi(y)-\phi(x)-S_k(y-x)\|
&\le(\zeta(S_k)-\|T(x)-S_k\|)\|x-y\|+\|T(x)-S_k\|\|x-y\|\cr
&\le\zeta(S_k)\|x-y\|\cr}$$

\noindent setiap kali $y\in D\cap B(x,\delta)$. Misalkan $q$,
$q'\in\Bbb Q^r$ sedemikian sehingga
$x\in[q,q']\subseteq B(x,\delta)$. Misalkan $n$ sedemikian sehingga
$q=q_n$, $q'=q'_n$ dan $k=m_n$. Maka $x\in H'_n$.\ \Qed

\medskip

{\bf (d)} Tuliskan

\Centerline{$C_n=\{x:x\in H_n,\,\lim_{\delta\downarrow 0}
\Bover{\mu^*(H_n\cap B(x,\delta))}{\mu B(x,\delta)}=1\}$.}

\noindent Maka $C_n\subseteq H'_n$.

\medskip

\Prf\ {\bf (i)} Ambil $x\in C_n$, dan tetapkan
$\tilde T=T(x)-S_{m_n}$. Saya harus menunjukkan bahwa
$\|\tilde T\|\le\zeta(S_{m_n})$. Ambil $\epsilon>0$. Misalkan
$\delta_0>0$ sedemikian sehingga

\Centerline{$\|\phi(y)-\phi(x)-T(x)(y-x)\|\le\epsilon\|y-x\|$}

\noindent setiap kali $y\in D$ dan $\|y-x\|\le\delta_0$.   Karena

\Centerline{$\|\phi(y)-\phi(x)-S_{m_n}(y-x)\|\le\zeta(S_{m_n})\|y-x\|$}

\noindent setiap kali $y\in H_n$, kita mempunyai

\Centerline{$\|\tilde T(y-x)\|\le(\epsilon+\zeta(S_{m_n}))\|y-x\|$}

\noindent setiap kali $y\in H_n$ dan $\|y-x\|\le\delta_0$.

\medskip

\quad{\bf (ii)} Berdasarkan 262L, terdapat $\delta_1>0$ sedemikian
sehingga $(1+2\epsilon)\delta_1\le\delta_0$ dan
$\rho(x+z,H_n)\le\epsilon\|z\|$ setiap kali
$0<\|z\|\le\delta_1$. Jadi jika $\|z\|\le\delta_1$, terdapat $y\in H_n$
sedemikian sehingga $\|x+z-y\|\le 2\epsilon\|z\|$. (Jika $z=0$, kita
dapat mengambil $y=x$.) Sekarang
$\|x-y\|\le (1+2\epsilon)\|z\|\le\delta_0$, sehingga

$$\eqalign{\|\tilde Tz\|
&\le\|\tilde T(y-x)\|+\|\tilde T(x+z-y)\|\cr
&\le(\epsilon+\zeta(S_{m_n}))\|y-x\|+\|\tilde T\|\|x+z-y\|\cr
&\le(\epsilon+\zeta(S_{m_n}))\|z\|
    +(\epsilon+\zeta(S_{m_n})+\|\tilde T\|)\|x+z-y\|\cr
&\le(\epsilon+\zeta(S_{m_n})+2\epsilon^2+2\epsilon\zeta(S_{m_n})
    +2\epsilon\|\tilde T\|)\|z\|.\cr}$$

\noindent Ini berlaku setiap kali $0<\|z\|\le\delta_1$.
Namun dengan mengalikan ketaksamaan ini dengan skalar positif yang sesuai,
kita melihat bahwa

\Centerline{$\|\tilde Tz\|
\le\bigl(\epsilon+\zeta(S_{m_n})+2\epsilon^2+2\epsilon\zeta(S_{m_n})
+2\epsilon\|\tilde T\|\bigr)\|z\|$}

\noindent untuk semua $z\in\BbbR^r$, dan

\Centerline{$\|\tilde T\|
\le\epsilon+\zeta(S_{m_n})+2\epsilon^2+2\epsilon\zeta(S_{m_n})
  +2\epsilon\|\tilde T\|$.}

\noindent Karena $\epsilon$ sebarang,
$\|\tilde T\|\le\zeta(S_{m_n})$, sebagaimana diklaim.\ \Qed

\medskip

{\bf (e)} Berdasarkan 261Da, $H_n\setminus C_n$ terabaikan untuk setiap
$n$, sehingga $H_n\setminus H'_n$ terabaikan, dan

\Centerline{$H'_n=D\cap(\overline{H}_n\setminus(H_n\setminus H'_n))$}

\noindent terukur relatif dalam $D$.   Tetapkan

\Centerline{$D_n=H'_n\setminus\bigcup_{k<n}H'_k$,
\quad$T_n=S_{m_n}$}

\noindent untuk setiap $n$; barisan-barisan ini memenuhi syarat.
}%end of proof of 262M


\leader{262N}{Akibat} Misalkan $\phi$ suatu fungsi dari himpunan bagian
$D$ dari $\Bbb R^r$ ke $\BbbR^s$, dan andaikan $\phi$ terdiferensialkan
relatif terhadap domainnya di setiap titik $D$. Maka $D$ dapat dinyatakan
sebagai gabungan barisan saling lepas $\sequencen{D_n}$ dari himpunan
bagian $D$ yang terukur relatif, sedemikian sehingga $\phi\restr D_n$
Lipschitz untuk setiap $n\in\Bbb N$.

\proof{ Dalam 262M, ambil $\zeta(T)=1$ untuk setiap $T\in A=M_{sr}$.
Jika $x$, $y\in D_n$, maka

$$\eqalign{\|\phi(x)-\phi(y)\|
&\le\|\phi(x)-\phi(y)-T_n(x-y)\|+\|T_n(x-y)\|\cr
&\le\|x-y\|+\|T_n\|\|x-y\|,\cr}$$

\noindent sehingga $\phi\restr D_n$ $(1+\|T_n\|)$-Lipschitz.
}%end of proof of 262N

\leader{262O}{Akibat} Andaikan $\phi$ suatu fungsi injektif dari himpunan
bagian terukur $D$ dari $\BbbR^r$ ke $\BbbR^r$, dan bahwa $\phi$
terdiferensialkan relatif terhadap domainnya di setiap titik $D$.

(a) Jika $A\subseteq D$ terabaikan, maka $\phi[A]$ terabaikan.

(b) Jika $E\subseteq D$ terukur, maka $\phi[E]$ terukur.

(c) Jika $D$ terukur dan $f$ suatu fungsi terukur yang didefinisikan pada
suatu himpunan bagian $\BbbR^r$, maka $f\phi^{-1}$ terukur.

(d) Jika $H\subseteq\BbbR^r$ dan $\phi^{-1}$ terdefinisi hampir di
mana-mana dalam $H$, dan jika $f$ suatu fungsi yang terdefinisi hampir di
mana-mana dalam $\Bbb R^r$, maka $f\phi^{-1}$ terdefinisi hampir di
mana-mana dalam $H$.

\proof{ Ambil barisan $\sequencen{D_n}$ seperti dalam 262N, dan terapkan
262E pada $\phi\restr D_n$ untuk setiap $n$.
}%end of proof of 262O

\leader{262P}{Akibat} Misalkan $\phi$ suatu fungsi dari suatu himpunan
bagian $D$ dari $\BbbR^r$ ke $\BbbR^s$, dan andaikan bahwa $\phi$
terdiferensialkan relatif terhadap domainnya, dengan suatu turunan
$T(x)$, pada setiap titik $x\in D$. Maka fungsi $x\mapsto T(x)$ terukur dalam
arti bahwa $\tau_{ij}:D\to\Bbb R$ terukur untuk semua $i\le s$ dan
$j\le r$, dengan $\tau_{ij}(x)$ koefisien ke-$(i,j)$ dari matriks $T(x)$
untuk semua $i$, $j$, dan $x$.

\proof{ Untuk setiap $k\in\Bbb N$, terapkan 262M dengan
$\zeta(T)=2^{-k}$ untuk setiap $T\in A=M_{sr}$, sehingga diperoleh
barisan $\sequencen{D_{kn}}$ dari himpunan bagian $D$ yang terukur
relatif, dan $\sequencen{T_{kn}}$ dalam $M_{sr}$. Misalkan
$\tau^{(kn)}_{ij}$ koefisien ke-$(i,j)$ dari $T_{kn}$. Maka kita
mempunyai fungsi $f_{ijk}:D\to\Bbb R$ yang didefinisikan dengan menetapkan

\Centerline{$f_{ijk}(x)=\tau^{(kn)}_{ij}$ jika $x\in D_{kn}$.}

\noindent Karena $D_{kn}$ terukur relatif, $f_{ijk}$ merupakan fungsi
terukur.   Untuk $x\in D_{kn}$,
\Centerline{$|\tau_{ij}(x)-f_{ijk}(x)|\le\|T(x)-T_{kn}\|\le 2^{-k}$,}

\noindent sehingga $|\tau_{ij}(x)-f_{ijk}(x)|\le 2^{-k}$ untuk setiap
$x\in D$, dan

\Centerline{$\tau_{ij}=\lim_{k\to\infty}f_{ijk}$}

\noindent terukur, sebagaimana diklaim.
}%end of proof of 262P

\leader{*262Q}{}\cmmnt{ Ini mengakhiri bagian yang esensial bagi sisa bab
ini. Namun menurut saya, hasil-hasil utama \S263 akan lebih mudah dipahami
jika Anda mengetahui fakta bahwa setiap fungsi Lipschitz terdiferensialkan
(relatif terhadap domainnya) hampir di mana-mana dalam domainnya. Saya
menyediakan beberapa halaman berikut untuk membuktikan fakta ini, yang
selain menarik pada dirinya sendiri juga merupakan latihan yang berguna.

\medskip

\noindent}{\bf Teorema Rademacher} Misalkan $\phi$ suatu fungsi Lipschitz
dari suatu himpunan bagian $\BbbR^r$ ke $\BbbR^s$. Maka $\phi$
terdiferensialkan relatif terhadap domainnya hampir di mana-mana dalam
domainnya.

\proof{{\bf (a)} Berdasarkan 262Ba dan 262Ib, cukup menangani kasus
$s=1$. Berdasarkan 262Bb, terdapat fungsi Lipschitz
$\tilde\phi:\BbbR^r\to\Bbb R$ yang memperluas $\phi$; kini $\phi$
terdiferensialkan relatif terhadap domainnya di setiap titik $\dom\phi$
tempat $\tilde\phi$ terdiferensialkan, sehingga cukup jika saya dapat
menunjukkan bahwa $\tilde\phi$ terdiferensialkan hampir di mana-mana.
Agar notasinya lebih nyaman dibaca, saya akan mengandaikan bahwa $\phi$
sendiri terdefinisi di seluruh $\BbbR^r$. Misalkan $\gamma$ suatu
konstanta Lipschitz bagi $\phi$.

Bukti berjalan dengan induksi pada $r$. Jika $r=1$, kita mempunyai fungsi
Lipschitz $\phi:\Bbb R\to\Bbb R$; kini $\phi$ kontinu mutlak dalam setiap
interval terbatas (262Bc), dan karena itu terdiferensialkan hampir di
mana-mana (225Cb). Dengan demikian induksi dimulai. Sisa bukti ditujukan
pada langkah induktif menuju $r>1$.

\medskip

{\bf (b)} Langkah pertama ialah menunjukkan bahwa semua turunan parsial
$\pd{\phi}{\xi_j}$ terdefinisi hampir di mana-mana dan terukur Borel.
\Prf\ Ambil $j\le r$. Untuk $q\in\Bbb Q\setminus\{0\}$ tetapkan

\Centerline{$\Delta_q(x)=\Bover{1}{q}(\phi(x+qe_j)-\phi(x))$,}

\noindent dengan menuliskan $e_j$ untuk vektor satuan ke-$j$ dari
$\BbbR^r$. Karena $\phi$ kontinu, $\Delta_q$ juga kontinu, sehingga
$\Delta_q$ merupakan fungsi terukur Borel untuk setiap $q$. Selanjutnya,
untuk sebarang $x\in\Bbb R^r$,

\Centerline{$D^+(x)=\limsup_{\delta\to 0}\Bover1{\delta}
(\phi(x+\delta e_j)-\phi(x))
=\lim_{n\to\infty}\sup_{q\in\Bbb Q,0<|q|\le 2^{-n}}\Delta_q(x)$,}

\noindent sehingga himpunan tempat $D^+(x)$ terdefinisi dalam $\Bbb R$
adalah himpunan Borel dan $D^+$ merupakan fungsi terukur Borel. Demikian
pula,

\Centerline{$D^-(x)=\liminf_{\delta\to 0}\Bover1{\delta}
(\phi(x+\delta e_j)-\phi(x))$}

\noindent merupakan fungsi terukur Borel dengan domain Borel.   Jadi

\Centerline{$E=\{x:\pd{\phi}{\xi_j}(x)\text{ ada dalam }\Bbb R\}
=\{x:D^+(x)=D^-(x)\in\Bbb R\}$}

\noindent merupakan himpunan Borel, dan $\pd{\phi}{\xi_j}$ merupakan
fungsi terukur Borel.

Di sisi lain, jika kita mengidentifikasi $\BbbR^r$ dengan
$\BbbR^J\times\Bbb R$, dengan mengambil $J$ sebagai
$\{1,\ldots,j-1,j+1,\ldots,r\}$, kita dapat memandang ukuran Lebesgue
$\mu$ pada $\BbbR^r$ sebagai produk ukuran Lebesgue $\mu_J$ pada
$\BbbR^J$ dengan ukuran Lebesgue $\mu_1$ pada $\Bbb R$ (251N). Kini,
untuk setiap $y\in\BbbR^J$, kita mempunyai fungsi
$\phi_y:\Bbb R\to\Bbb R$ yang didefinisikan dengan menuliskan

\Centerline{$\phi_y(\sigma)=\phi(y,\sigma)$,}

\noindent dan $E$ menjadi

\Centerline{$\{(y,\sigma):\phi_y'(\sigma)$ terdefinisi$\}$,}

\noindent sehingga semua penampang

\Centerline{$\{\sigma:(y,\sigma)\in E\}$}

\noindent merupakan himpunan bagian koterabaikan dari $\Bbb R$, sebab
setiap $\phi_y$ Lipschitz, dan karena itu terdiferensialkan hampir di
mana-mana, sebagaimana dicatat dalam bagian (a) bukti. Karena kita tahu
bahwa $E$ terukur, himpunan tersebut harus koterabaikan berdasarkan teorema
Fubini (terapkan 252D atau 252F pada komplemen $E$). Dengan demikian
$\pd{\phi}{\xi_j}$ terdefinisi hampir di mana-mana, sebagaimana
diklaim.\ \Qed

Tuliskan

\Centerline{$H=\{x:x\in\BbbR^r,\,\pd{\phi}{\xi_j}(x)
\text{ ada untuk setiap }j\le r\}$,}

\noindent sehingga $H$ merupakan himpunan Borel koterabaikan dalam
$\BbbR^r$.

\medskip

{\bf (c)} Untuk sisa bukti ini, saya tetapkan identifikasi alami
$\BbbR^r$ dengan $\BbbR^{r-1}\times\Bbb R$, dengan mengidentifikasi
$(\xi_1,\ldots,\xi_r)$ dengan $((\xi_1,\ldots,\xi_{r-1}),\xi_r)$.
Untuk $x\in H$, misalkan $T(x)$ matriks $1\times r$
$(\pd{\phi}{\xi_1}(x),\ldots,\pd{\phi}{\xi_r}(x))$.

\medskip

{\bf (d)} Tetapkan

\Centerline{$H_1=\{x:x\in H,\,\lim_{u\to\tbf{0}\text{ dalam }\BbbR^{r-1}}
\Bover{|\phi(x+(u,0))-\phi(x)-T(x)(u,0)|}{\|u\|}=0\}$.}

\noindent Saya nyatakan bahwa $H_1$ koterabaikan dalam $\BbbR^r$.
\Prf\ Ini sebenarnya gagasan yang sama seperti dalam (b). Untuk
$x\in H$, $x\in H_1$ jika dan hanya jika

\inset{untuk setiap $\epsilon>0$ terdapat $\delta>0$ sedemikian sehingga

\Centerline{$|\phi(x+(u,0))-\phi(x)-T(x)(u,0)|\le\epsilon\|u\|$}

\noindent setiap kali $\|u\|\le\delta$,}

\noindent yaitu, jika dan hanya jika

\inset{untuk setiap $m\in\Bbb N$ terdapat $n\in\Bbb N$ sedemikian
sehingga

\Centerline{$|\phi(x+(u,0))-\phi(x)-T(x)(u,0)|\le 2^{-m}\|u\|$}

\noindent setiap kali $u\in\Bbb Q^{r-1}$ dan $\|u\|\le 2^{-n}$.}

\noindent Namun, untuk setiap $m\in\Bbb N$ dan $u\in\Bbb Q^{r-1}$ tertentu,
himpunan

\Centerline{$\{x:|\phi(x+(u,0))-\phi(x)-T(x)(u,0)|\le 2^{-m}\|u\|\}$}

\noindent terukur, bahkan Borel, sebab semua fungsi
$x\mapsto\phi(x+(u,0))$, $x\mapsto\phi(x)$, $x\mapsto T(x)(u,0)$ terukur
Borel. Jadi $H_1$ berbentuk

\Centerline{$\bigcap_{m\in\Bbb N}\bigcup_{n\in\Bbb N}
\bigcap_{u\in\Bbb Q^{r-1},\|u\|\le 2^{-n}}E_{mnu}$}

\noindent dengan setiap $E_{mnu}$ suatu himpunan terukur, sehingga $H_1$
terukur.

Namun sekarang perhatikan bahwa untuk sebarang $\sigma\in\Bbb R$, fungsi

\Centerline{$v\mapsto\phi_{\sigma}(v)=\phi(v,\sigma):
\BbbR^{r-1}\to\Bbb R$}

\noindent Lipschitz, dan karena itu (berdasarkan hipotesis induktif)
terdiferensialkan hampir di mana-mana dalam $\BbbR^{r-1}$; dan
$(v,\sigma)\in H_1$ jika dan hanya jika $(v,\sigma)\in H$ dan
$\phi_{\sigma}'(v)$ terdefinisi. Akibatnya
$\{v:(v,\sigma)\in H_1\}$ koterabaikan setiap kali
$\{v:(v,\sigma)\in H\}$ koterabaikan, yaitu untuk hampir setiap
$\sigma\in\Bbb R$; sehingga $H_1$, karena terukur, harus koterabaikan.\ \Qed

\medskip

{\bf (e)} Sekarang, untuk $q$, $q'\in\Bbb Q$ dan $n\in\Bbb N$, tetapkan

\Centerline{$F(q,q',n)
=\{x:x\in\BbbR^r$,
  $q\le\Bover{\phi(x+(\tbf{0},\eta))-\phi(x)}{\eta}\le q'$
  setiap kali $0<|\eta|\le 2^{-n}\}$,}

\Centerline{$F_*(q,q',n)
=\{x:x\in F(q,q',n),\,\lim_{\delta\downarrow 0}
 \Bover{\mu^*(F(q,q',n)\cap B(x,\delta))}{\mu B(x,\delta)}=1\}$.}

\noindent Berdasarkan 261Da, $F(q,q',n)\setminus F_*(q,q',n)$ terabaikan
untuk semua $q$, $q'$, $n$, sehingga

\Centerline{$H_2
=H_1\setminus\bigcup_{q,q'\in\Bbb Q,n\in\Bbb N}
  (F(q,q',n)\setminus F_*(q,q',n))$}

\noindent koterabaikan.

\medskip

{\bf (f)} Saya nyatakan bahwa $\phi$ terdiferensialkan di setiap titik
$H_2$. \Prf\ Ambil $x=(u,\sigma)\in H_2$. Maka
$\alpha=\Pd{\phi}{\xi_r}(x)$ dan $T=T(x)$ terdefinisi. Misalkan $\gamma$
suatu konstanta Lipschitz bagi $\phi$.

Ambil $\epsilon>0$; ambil $q$, $q'\in\Bbb Q$ sedemikian sehingga
$\alpha-\epsilon\le q<\alpha<q'\le\alpha+\epsilon$. Harus ada
$n\in\Bbb N$ sedemikian sehingga $x\in F(q,q',n)$; akibatnya
$x\in F_*(q,q',n)$ berdasarkan definisi $H_2$. Berdasarkan 262L, terdapat
$\delta_0>0$ sedemikian sehingga
$\rho(x+z,F(q,q',n))\le\epsilon\|z\|$ setiap kali $\|z\|\le\delta_0$.
Selanjutnya, terdapat $\delta_1>0$ sedemikian sehingga
$|\phi(x+(v,0))-\phi(x)-T(v,0)|\le\epsilon\|v\|$ setiap kali
$v\in\BbbR^{r-1}$ dan $\|v\|\le\delta_1$. Tetapkan

\Centerline{$\delta=\min(\delta_0,\delta_1,2^{-n})/(1+2\epsilon)>0$.}

Andaikan $z=(v,\tau)\in\BbbR^r$ dan $\|z\|\le\delta$. Karena
$\|z\|\le\delta_0$, terdapat $x'=(u',\sigma')\in F(q,q',n)$ sedemikian
sehingga $\|x+z-x'\|\le2\epsilon\|z\|$; tetapkan $x^*=(u',\sigma)$.
Sekarang

\Centerline{$\max(\|u-u'\|,|\sigma-\sigma'|)\le
\|x-x'\|\le(1+2\epsilon)\|z\|\le\min(\delta_1,2^{-n})$}

\noindent dan $x^*=x+(u'-u,0)$, sehingga

\Centerline{$|\phi(x^*)-\phi(x)-T(x^*-x)|
\le\epsilon\|u'-u\|\le\epsilon(1+2\epsilon)\|z\|$.}

\noindent Tetapi juga

\Centerline{$|\phi(x')-\phi(x^*)-T(x'-x^*)|
=|\phi(x')-\phi(x^*)-\alpha(\sigma'-\sigma)|
\le\epsilon|\sigma'-\sigma|\le\epsilon(1+ 2\epsilon)\|z\|$,}

\noindent sebab $x'\in F(q,q',n)$ dan
$|\sigma-\sigma'|\le 2^{-n}$, sehingga (jika $x'\ne x^*$)

\Centerline{$\alpha-\epsilon\le q
\le\Bover{\phi(x^*)-\phi(x')}{\sigma-\sigma'}\le q'\le\alpha+\epsilon$}

\noindent dan

\Centerline{$\bigl|\Bover{\phi(x')-\phi(x^*)}
{\sigma'-\sigma}-\alpha\bigr|\le\epsilon$.}
\noindent Akhirnya,

\Centerline{$|\phi(x+z)-\phi(x')|
\le\gamma\|x+z-x'\|\le2\gamma\epsilon\|z\|$,}

\Centerline{$|Tz-T(x'-x)|\le\|T\|\|x+z-x'\|\le2\epsilon\|T\|\|z\|$.}

Dengan menggabungkan semuanya,

$$\eqalign{|\phi(x+z)-\phi x-Tz|
&\le|\phi(x+z)-\phi(x')|
+|T(x'-x)-Tz|\cr
&\qquad
+|\phi(x')-\phi(x^*)-T(x'-x^*)|
+|\phi(x^*)-\phi(x)-T(x^*-x)|\cr
&\le2\gamma\epsilon\|z\|
+2\epsilon\|T\|\|z\|
+\epsilon(1+2\epsilon)\|z\|
+\epsilon(1+2\epsilon)\|z\|\cr
&=\epsilon(2\gamma+2\|T\|+2+4\epsilon)\|z\|.\cr}$$

\noindent Ini berlaku setiap kali $\|z\|\le\delta$. Karena $\epsilon$
sebarang, $\phi$ terdiferensialkan di $x$.\ \Qed

Dengan demikian, $\{x:\phi$ terdiferensialkan di $x\}$ memuat $H_2$ dan
koterabaikan; induksi pun berlanjut.
}%end of proof of 262Q

\exercises{
\leader{262X}{Latihan dasar (a)} Misalkan $\phi$ dan $\psi$ fungsi
Lipschitz dari himpunan-himpunan bagian $\BbbR^r$ ke $\BbbR^s$.
Tunjukkan bahwa $\phi+\psi$ merupakan fungsi Lipschitz dari
$\dom\phi\cap\dom\psi$ ke
$\Bbb R^s$.
%262A

\header{262Xb}{\bf (b)} Misalkan $\phi$ fungsi Lipschitz dari suatu
himpunan bagian $\BbbR^r$ ke $\BbbR^s$, dan $c\in\Bbb R$. Tunjukkan
bahwa $c\phi$ merupakan fungsi Lipschitz.
%262A

\header{262Xc}{\bf (c)} Andaikan $\phi:D\to\BbbR^s$ dan
$\psi:E\to\BbbR^q$ fungsi-fungsi Lipschitz, dengan
$D\subseteq\BbbR^r$ dan $E\subseteq\BbbR^s$. Tunjukkan bahwa komposisi
$\psi\phi:D\cap\phi^{-1}[E]\to\BbbR^q$ bersifat Lipschitz.
%262A

\header{262Xd}{\bf (d)} Andaikan $\phi$, $\psi$ fungsi-fungsi dari
himpunan-himpunan bagian $\BbbR^r$ ke $\BbbR^s$, dan andaikan
$x\in\dom\phi\cap\dom\psi$ sedemikian sehingga masing-masing fungsi
terdiferensialkan relatif terhadap domainnya di $x$, dengan turunan
$S$, $T$ di sana. Tunjukkan bahwa $\phi+\psi$ terdiferensialkan relatif
terhadap domainnya di $x$, dan bahwa $S+T$ merupakan suatu turunan dari
$\phi+\psi$ di $x$.
%262F

\header{262Xe}{\bf (e)} Andaikan $\phi$ suatu fungsi dari suatu himpunan
bagian $\BbbR^r$ ke $\BbbR^s$, dan terdiferensialkan relatif terhadap
domainnya di $x\in\dom\phi$. Tunjukkan bahwa $c\phi$ terdiferensialkan
relatif terhadap domainnya di $x$ untuk setiap $c\in\Bbb R$.
%262F

\sqheader 262Xf Andaikan $\phi:D\to\BbbR^s$ dan
$\psi:E\to\Bbb R^q$ fungsi-fungsi, dengan
$D\subseteq\BbbR^r$ dan $E\subseteq\BbbR^s$; andaikan $\phi$
terdiferensialkan relatif terhadap domainnya di
$x\in D\cap\phi^{-1}[E]$, dengan matriks $s\times r$, yaitu $T$, sebagai
suatu turunan di sana, dan $\psi$ terdiferensialkan relatif terhadap
domainnya di $\phi(x)$, dengan matriks $q\times s$, yaitu $S$, sebagai
suatu turunan di sana. Tunjukkan bahwa komposisi $\psi\phi$
terdiferensialkan relatif terhadap domainnya di $x$, dan bahwa matriks
$q\times r$, yaitu $ST$, merupakan suatu turunan dari $\psi\phi$ di $x$.
%262F

\sqheader 262Xg Misalkan $\phi:\BbbR^r\to\BbbR^s$ suatu operator linear,
dengan matriks terkait $T$. Tunjukkan bahwa $\phi$ terdiferensialkan
di mana-mana, dengan $\phi'(x)=T$ untuk setiap $x$.
%262F

\spheader 262Xh Misalkan $G\subseteq\BbbR^r$ suatu himpunan terbuka
konveks, dan $\phi:G\to\BbbR^s$ suatu fungsi sedemikian sehingga semua
turunan parsial $\pd{\phi_i}{\xi_j}$ terdefinisi di mana-mana dalam $G$.
Tunjukkan bahwa $\phi$ Lipschitz jika dan hanya jika semua turunan
parsialnya terbatas pada $G$.
%262I

\spheader 262Xi Misalkan $\phi:\BbbR^r\to\BbbR^s$ suatu fungsi.
Tunjukkan bahwa $\phi$ terdiferensialkan di $x\in\BbbR^r$ jika dan hanya
jika untuk setiap $m\in\Bbb N$ terdapat $n\in\Bbb N$ dan matriks
$s\times r$, yaitu $T$, dengan koefisien rasional sedemikian sehingga
$\|\phi(y)-\phi(x)-T(y-x)\|\le 2^{-m}\|y-x\|$ setiap kali
$\|y-x\|\le 2^{-n}$.
%262I

\spheader 262Xj Andaikan $f$ suatu fungsi bernilai real yang terintegralkan
atas $\BbbR^r$, dan $g:\BbbR^r\to\Bbb R$ suatu fungsi terdiferensialkan
dan terbatas sedemikian sehingga turunan parsial $\Pd{g}{\xi_j}$ terbatas,
dengan $j\le r$. Misalkan $f*g$ konvolusi $f$ dan $g$ (255E, 255L).
Tunjukkan bahwa $\Pd{}{\xi_j}(f*g)$ terdefinisi di mana-mana dan sama
dengan $f*\Pd{g}{\xi_j}$. \Hint{255Xd.}
%262+

\spheader 262Xk Misalkan $(X,\Sigma,\mu)$ suatu ruang ukuran,
$G\subseteq\BbbR^r$ suatu himpunan terbuka, dan
$f:X\times G\to\Bbb R$ suatu fungsi. Andaikan bahwa

\inset{(i) untuk setiap $x\in X$, fungsi
$t\mapsto f(x,t):G\to\Bbb R$ terdiferensialkan;

(ii) terdapat fungsi terintegralkan $g$ pada $X$ sedemikian sehingga
$|\pd{f}{\tau_j}(x,t)|\le g(x)$ setiap kali $x\in X$, $t\in G$ dan
$j\le r$;

(iii) $\int|f(x,t)|\mu(dx)$ ada dalam $\Bbb R$ untuk setiap $t\in G$.}

\noindent Tunjukkan bahwa fungsi
$t\mapsto\int f(x,t)\mu(dx):G\to\Bbb R$ terdiferensialkan.
\Hint{tunjukkan terlebih dahulu bahwa, untuk suatu $M$ yang sesuai,
$|f(x,t)-f(x,t')|\le M|g(x)|\|t-t'\|$ untuk setiap $t$, $t'\in G$ dan
$x\in X$.}
%262+

\spheader 262Xl\dvAnew{2015} Misalkan $f:[a,b]\to\Bbb R$ suatu fungsi
kontinu mutlak, dengan $a\le b$, dan $g:f[\,[a,b]\,]\to\Bbb R$ suatu
fungsi Lipschitz. Tunjukkan bahwa $gf$ kontinu mutlak.
%262Bc out of order query

\leader{262Y}{Latihan lanjutan (a)}
%\spheader 262Ya
Misalkan $L$ ruang semua fungsi Lipschitz dari $\BbbR^r$ ke $\BbbR^s$,
dan untuk $\phi\in L$ tetapkan

\Centerline{$\|\phi\|
=\|\phi(0)\|+\min\{\gamma:\gamma\in\coint{0,\infty}$,
$\|\phi(y)-\phi(x)\|\le\gamma\|y-x\|$ untuk setiap $x$, $y\in\BbbR^r\}$.}

\noindent Tunjukkan bahwa $(L,\|\,\,\|)$ merupakan ruang Banach.
%262A

\spheader 262Yb Tunjukkan bahwa jika
$T=\langle\tau_{ij}\rangle_{i\le s,j\le r}$ suatu matriks $s\times r$,
maka norma operator $\|T\|$, sebagaimana didefinisikan dalam 262H, paling
besar
$\sqrt{\sum_{i=1}^s\sum_{j=1}^r\tau_{ij}^2}$.
%262H

\spheader 262Yc Misalkan $\phi:D\to\Bbb R$ sebarang fungsi, dengan
$D\subseteq\BbbR^r$. Tunjukkan bahwa
$H=\{x:x\in D,\,\phi$ terdiferensialkan relatif terhadap domainnya di
$x\}$ terukur relatif dalam $D$, dan bahwa $\pd{\phi}{\xi_j}\restr H$
terukur untuk setiap $j\le r$.
%262I

\spheader 262Yd\dvAnew{2015}
Misalkan $\phi:D\to\Bbb R$ suatu fungsi, dengan
$D\subseteq\BbbR^r$.
(i) Tunjukkan bahwa jika $\phi$ terukur, maka semua turunan parsialnya
terukur.
(ii) Tunjukkan bahwa jika $\phi$ terukur Borel, maka semua turunan
parsialnya terukur Borel.
%262Yc 262I 225Yg

\spheader 262Ye Suatu fungsi $\phi:\BbbR^r\to\Bbb R$ disebut
{\bf mulus} jika semua turunan parsialnya
$\Bover{\partial\ldots\partial\phi}
{\partial\xi_i\partial\xi_j\ldots\partial\xi_l}$ terdefinisi di mana-mana
dalam $\BbbR^r$ dan kontinu. Tunjukkan bahwa jika $f$ terintegralkan atas
$\BbbR^r$ dan $\phi:\BbbR^r\to\Bbb R$ mulus serta mempunyai dukungan
terbatas, maka konvolusi $f*\phi$ mulus. \Hint{262Xj, 262Xk.}
%262Xk 262+

\spheader 262Yf Untuk $\delta>0$, tetapkan
$\tilde\phi_{\delta}(x) =e^{-1/(\delta^2-\|x\|^2)}$ jika $\|x\|<\delta$,
$0$ jika $\|x\|\ge\delta$; tetapkan
$\alpha_{\delta}=\int\tilde\phi_{\delta}(x)dx$,
$\phi_{\delta}(x)=\alpha_{\delta}^{-1}\tilde\phi_{\delta}(x)$ untuk setiap
$x$. (i) Tunjukkan bahwa $\phi_{\delta}:\BbbR^r\to\Bbb R$ mulus dan
mempunyai dukungan terbatas. (ii) Tunjukkan bahwa jika $f$ terintegralkan
atas $\BbbR^r$, maka
$\lim_{\delta\downarrow 0}\int|f(x)-(f*\phi_{\delta})(x)|dx=0$.
\Hint{mulailah dengan fungsi-fungsi kontinu $f$ yang mempunyai dukungan
terbatas, dan gunakan 242O.}
%262Ye 262+

\spheader 262Yg Tunjukkan bahwa jika $f$ terintegralkan atas $\BbbR^r$
dan $\epsilon>0$, terdapat fungsi mulus $h$ dengan dukungan terbatas
sedemikian sehingga $\int|f-h|\le\epsilon$.
\Hint{{\it salah satu}: reduksikan ke kasus ketika $f$ mempunyai dukungan
terbatas dan gunakan 262Yf; {\it atau} sesuaikan metode 242Xi.}
%262Yf 262+

\spheader 262Yh Andaikan $f$ suatu fungsi real yang terintegralkan atas
setiap himpunan bagian terbatas $\BbbR^r$. (i) Tunjukkan bahwa
$f\times\phi$ terintegralkan setiap kali $\phi:\BbbR^r\to\Bbb R$ suatu
fungsi mulus dengan dukungan terbatas. (ii) Tunjukkan bahwa jika
$\int f\times\phi=0$ untuk setiap fungsi mulus dengan dukungan terbatas,
maka $f=0$ hampir di mana-mana. ({\it Petunjuk\/}: tunjukkan bahwa
$\int_{B(x,\delta)}f=0$ untuk setiap $x\in\BbbR^r$ dan $\delta>0$, dan
gunakan 261C. Sebagai alternatif, tunjukkan terlebih dahulu bahwa
$\int_Ef=0$ untuk $E=[b,c]$, kemudian untuk himpunan terbuka $E$, lalu
untuk sebarang himpunan terukur $E$.)
%262Ye 262+

\spheader 262Yi Misalkan $f$ terintegralkan atas $\BbbR^r$, dan untuk
$\delta>0$ misalkan $\phi_{\delta}:\BbbR^r\to\Bbb R$ fungsi dari 262Yf.
Tunjukkan bahwa $\lim_{\delta\downarrow 0}(f*\phi_{\delta})(x)=f(x)$
untuk setiap $x$ dalam himpunan Lebesgue $f$. ({\it Petunjuk\/}: 261Ye.)
%262Yf 262+

}%end of exercises

\endnotes{
\Notesheader{262} Penekanan bagian ini ternyata terletak pada hubungan
antara konsep `fungsi Lipschitz' dan `fungsi terdiferensialkan'. Salah
satu pesona analisis real klasik ialah munculnya hubungan sedemikian erat
antara konsep-konsep yang termasuk dalam kategori berbeda. `Fungsi
Lipschitz' jelas termasuk dalam teori ruang metrik (saya akan kembali ke
hal ini dalam \S264), sedangkan `fungsi terdiferensialkan' termasuk dalam
teori manifold terdiferensialkan, yang berada di luar cakupan jilid ini.
Saya telah menuliskan bagian ini dengan saksama untuk berjaga-jaga jika
ada pembaca yang sejauh ini belum mempelajari teori pemetaan
terdiferensialkan antar-ruang Euklides multidimensi; tetapi ini juga
memberi saya kesempatan untuk menguraikan gagasan `fungsi yang
terdiferensialkan relatif terhadap domainnya', yang memungkinkan kita
menangani dengan lancar berbagai masalah yang muncul pada batas dalam
bagian berikutnya. Kesulitan yang saya perhatikan pertama-tama muncul
pada fungsi seperti transformasi koordinat polar

\Centerline{$(\rho,\theta)\mapsto (\rho\cos\theta,\rho\sin\theta):
\{(0,0)\}\cup(\ooint{0,\infty}\times\ocint{-\pi,\pi})\to\BbbR^2$.}

\noindent Untuk menjadikannya bijeksi, kita harus melakukan sesuatu yang
agak sebarang, dan domain transformasi tersebut tidak mungkin berupa
himpunan terbuka. Berdasarkan definisi yang saya gunakan, fungsi ini
terdiferensialkan relatif terhadap domainnya di setiap titik domainnya,
dan kita dapat menerapkan hasil seperti 262O tanpa hambatan. Perhatikan
bahwa dalam kasus ini titik-titik domain yang bukan titik interior
membentuk himpunan terabaikan
$\{(0,0)\}\cup(\ooint{0,\infty}\times\{\pi\})$, sehingga kita dapat
berharap untuk mengabaikannya; dan bagi sebagian besar transformasi
geometris langsung yang menjadi sasaran penerapan teori ini, penghilangan
himpunan-himpunan terabaikan secara bijaksana akan mereduksi masalah ke
kasus fungsi yang terdiferensialkan dalam arti biasa dengan domain terbuka. Namun,
meskipun teori domain terbuka akan menangani sebagian besar contoh yang
paling penting, ada bahaya bahwa Anda akan memiliki kesalahpahaman nyata
mengenai cakupan metode-metode ini.

Hakikat keterdiferensialan ialah bahwa suatu fungsi terdiferensialkan
$\phi$ dapat dihampiri, di dekat sebarang titik tertentu dalam domainnya,
oleh fungsi afin. Gagasan 262M ialah menguraikan suatu metode yang efektif
untuk cakupan luas guna membagi $D=\dom\phi$ menjadi terhitung banyak bagian,
dan pada setiap bagian itu $\phi$ terkendali dengan baik. Metode ini akan
diterapkan dalam \S\S263 dan 265 untuk menyelidiki ukuran $\phi[D]$;
tetapi kita sudah mempunyai beberapa akibat langsung (262N-262P).

Saya telah memberikan sejumlah hasil yang menunjukkan bahwa (berdasarkan
definisi-definisi yang saya pilih) suatu turunan dapat diharapkan memiliki
sekurang-kurangnya beberapa sifat `deskriptif' yang juga dimiliki fungsi
asalnya; lihat 222Yd, 225J, 225Yg, 262Yc, 262Yd. Untuk turunan parsial,
terdapat komplikasi mengenai domainnya (419Yd, 431Yd) yang tidak muncul
pada turunan penuh (225J, 262Yc).
}%end of notes

\discrpage
